\documentclass[11pt,a4paper]{article}
\usepackage[margin=1in]{geometry}
\usepackage{hyperref}
\usepackage{enumitem}
\usepackage{graphicx}
\usepackage{xcolor}

% -----------------------------------------------------
% -----------------------------------------------------

% Variable: Lab Instructor
\makeatletter \newcommand{\BvrSetLabInstructor}[1]{%
  \gdef\Bvr@LabInstructor{#1}}
% default
\newcommand{\Bvr@LabInstructor}{Dr. Nisha Thakur}
% public accessor
\newcommand{\BvrLabInstructorValue}{\Bvr@LabInstructor}
\makeatother

% Add subtitle
\usepackage{titling}
% -- Set title bold
\pretitle{\begin{center}\bfseries\fontsize{18bp}{18bp}\selectfont}
\makeatletter
\newcommand{\BvrSubtitle}[1]{%
  \posttitle{%
    \par\end{center}
    \begin{center}\large#1\end{center}
    \vskip 0.5em}%
}
\makeatother

% -----------------------------------------------------

\title{POTLI: On demand Luggage storage network}

\BvrSetLabInstructor{Nisha Thakur}

\BvrSubtitle{%
  Project Proposal for UCS503P  \\
  Submitted to: \BvrLabInstructorValue \\ \bfseries
  Thapar Institute of Engineering and Technology% 
}

\author{
  Satyam Tiwari, CSED%
  \thanks{Roll No: \texttt{1024030088}, %
    Email: \texttt{stiwari1\_be24@thapar.edu}}
  \and
  Ishant Mehndiratta, CSED%
  \thanks{Roll No: \texttt{1024030525}, %
    Email: \texttt{imehndiratta\_be24@thapar.edu}}
  \and
  Anshaj, CSED%
  \thanks{Roll No: \texttt{1024030494}, %
    Email: \texttt{aanshaj\_be24@thapar.edu}}
  \and
  Aayush Bindal, CSED%
  \thanks{Roll No: \texttt{1024030498}, %
    Email: \texttt{abindal1\_be24@thapar.edu}}
}

\date{\today}

\begin{document}
\maketitle

\section{Higher-order goal}
This project aims to make short-distance urban travel a little less stressful
by giving people a safe, nearby place to leave their bags whenever they don't
want to lug them around -- whether they're killing time before a hotel
check-in, fitting in some sightseeing between train connections, or just
running errands across an unfamiliar city. While the broader theme is about
smoothing out everyday urban travel and opening up a genuine extra income
stream for small local businesses, the immediate engineering objective is to
translate this idea into a measurable, deployable software solution.

\section{Time-to-value}
\begin{itemize}[leftmargin=0.7cm]
    \item We will deliver meaningful increments quickly by prototyping the core booking-and-handoff workflow and validating it with a small set of real partner shops within a short iteration window.
    \item We will prioritize fast feedback over perfection by using weekly iterations (and targeting CI-backed automation early) to reduce release friction.
    \item Success is defined by what can be delivered and measured in the first iteration, then improved based on partner and traveler feedback.
\end{itemize}

\section{Problem Statement}
Travelers and short-term visitors in Indian cities -- tourists between
check-in and check-out, students home for the day, professionals in town for
a single meeting -- regularly need somewhere safe to leave a bag for a few
hours, but have almost nowhere reliable to turn. Common pain points include:
\begin{itemize}[leftmargin=0.7cm]
    \item \textbf{Scarcity:} official lockers exist only at a handful of railway stations and airports, and are frequently full, expensive, or out of service.
    \item \textbf{Low discoverability:} even when a nearby shop or business might be willing to hold a bag, travelers have no easy way to find or trust that option.
    \item \textbf{Trust and safety:} leaving belongings with a stranger, with no verification, receipt, or accountability, is understandably uncomfortable.
    \item \textbf{Untapped local income:} small shop and business owners often have a spare corner, shelf, or storeroom sitting idle, with no simple way to earn from it.
\end{itemize}

\textbf{Why this matters now (contextual relevance):} as more people travel for
work, tourism, and short visits, and as tier-2/3 cities see rising footfall
without matching storage infrastructure, even a lightweight coordination
layer between travelers and local businesses can noticeably reduce this
everyday friction.

\section{Proposed Solution}
\subsection{Overview}
Build a \textbf{web-based} ``Storage-on-Demand'' network with the following components:
\begin{itemize}[leftmargin=0.7cm]
    \item \textbf{Traveler discovery \& booking portal:} search nearby partner shops, see capacity/price/hours, and book a slot in advance or on the spot.
    \item \textbf{Partner onboarding portal:} shop and business owners list their available space, set capacity, pricing, and working hours, and manage incoming bookings.
    \item \textbf{Booking \& handoff workflow:} QR-code based check-in and check-out so both sides can verify exactly what was dropped off and what was picked up.
    \item \textbf{Trust \& safety layer:} verified partner listings, photo confirmation at drop-off, and post-pickup ratings.
    \item \textbf{Payments \& payouts:} the traveler pays through the platform; the partner receives their share automatically once the booking is completed.
    \item \textbf{Admin/ops dashboard:} city-wise coverage, partner performance, and revenue summaries.
\end{itemize}

\subsection{Core workflow}
\begin{enumerate}[leftmargin=0.7cm]
    \item Traveler searches for storage near a location and picks a partner and a time window.
    \item System generates a booking ID and a QR code for that booking.
    \item Partner scans/confirms the bag at drop-off; the bag is tagged with the booking ID.
    \item Traveler returns within the window and the partner scans/confirms pickup.
    \item Payment is settled and the partner's payout is credited.
\end{enumerate}

\subsection{Operational constraints for an educational, tier-2/3 setting}
\begin{itemize}[leftmargin=0.7cm]
    \item Keep the solution usable on low-to-mid range devices, since many partner shop owners will not have high-end smartphones.
    \item Keep partner onboarding extremely simple -- most partners will be small shopkeepers, not tech-savvy businesses.
    \item Avoid dependence on advanced research algorithms; focus on workflow correctness, trust, and reliable handoffs.
    \item Design for deployability using common web hosting and managed databases.
\end{itemize}

\section{Solution Approach}
This is an engineering course project, so we focus on scalable coordination and trust-related workflow aspects rather than novel algorithm research.

\subsection{Partner listing \& duplicate-signup checks}
Duplicate/quality checks on partner listings will be heuristic-based:
\begin{itemize}[leftmargin=0.7cm]
    \item Normalize partner details (business name, locality/address, phone number).
    \item Use a simple similarity measure (e.g., name + locality token overlap) to flag likely duplicate signups, without claiming any state-of-the-art matching technique.
    \item Show flagged listings to an admin for review; never auto-approve or auto-reject a partner without a human check.
\end{itemize}

\subsection{Web architecture}
A typical 3-tier web architecture:
\begin{itemize}[leftmargin=0.7cm]
    \item \textbf{Frontend:} responsive UI for two user types -- travelers (search, book, track) and partners (list space, manage bookings).
    \item \textbf{Backend API:} authentication, location-based search, booking and handoff state transitions, payments.
    \item \textbf{Database:} users, partners, listings, bookings, handoff logs, and transactions.
\end{itemize}

\subsection{CI/CD and fast delivery enabler}
CI/CD will be used to:
\begin{itemize}[leftmargin=0.7cm]
    \item Automatically build and run tests on every change.
    \item Produce repeatable deployment artifacts to staging.
    \item Enable safe and frequent releases of incremental features (e.g., shipping QR-based handoff before layering on payments).
\end{itemize}

\section{Evaluation Criterion}
We define success using measurable metrics tied to the core booking-and-handoff workflow.

\subsection{Primary evaluation metric}
\textbf{Booking Fulfillment Time (BFT):} time between a traveler starting a search and receiving a confirmed booking at a partner location.
\begin{itemize}[leftmargin=0.7cm]
    \item Target: median BFT $\le$ 3 minutes during the pilot window.
    \item Attribution: measured from database timestamps (search/request vs. confirmed booking).
\end{itemize}

\subsection{Secondary evaluation metrics}
\begin{itemize}[leftmargin=0.7cm]
    \item \textbf{Handoff success rate:} percentage of bookings where both drop-off and pickup are confirmed without disputes.
    \item \textbf{Partner utilization:} percentage of a partner's listed capacity actually booked per week.
    \item \textbf{Partner payout turnaround:} time between booking completion and payout being credited.
    \item \textbf{Traveler satisfaction:} post-pickup 1--5 rating on ease and trust.
    \item \textbf{Reliability:} platform availability for search/booking during business hours (e.g., $\ge$ 99\% uptime in pilot).
\end{itemize}

\subsection{Pilot validation plan}
\begin{itemize}[leftmargin=0.7cm]
    \item Recruit 8--12 partner shops/businesses in a chosen locality (e.g., near a railway station or a busy market) and 15--20 traveler participants.
    \item Compare metrics against a baseline of no formal storage option (carrying bags around or informal arrangements), collected via short interviews before the pilot.
\end{itemize}

\section{Scalability}
The proposition should scale in both theoretical and practical senses.

\begin{enumerate}[leftmargin=0.7cm]
    \item \textbf{Scalable (theoretically):} The system should support growth in listings, bookings, and concurrent searches by:
    \begin{itemize}[leftmargin=0.7cm]
        \item decoupling components (frontend/backend),
        \item using geo-indexed queries for ``storage near me'' search and pagination for results,
        \item using stateless API design.
    \end{itemize}
    
    \item \textbf{Operationally deployable (practically):} The system should run on common web hosting:
    \begin{itemize}[leftmargin=0.7cm]
        \item managed database,
        \item containerized or platform-hosted deployment,
        \item CI/CD pipeline for staging/production.
    \end{itemize}
\end{enumerate}

\section{Engine availability heuristic}
A mature set of ``engines'' (tooling/services) is available to implement this as a web-based project:
\begin{itemize}[leftmargin=0.7cm]
    \item Web framework for REST APIs and responsive pages.
    \item Managed database with geospatial/location query support.
    \item Authentication approach (token-based or session-based) for both traveler and partner accounts.
    \item QR-code generation/scanning library for handoff verification.
    \item Payment gateway integration (e.g., UPI-based) for bookings and partner payouts.
    \item Notification channel (in-app, optionally SMS/WhatsApp) for booking confirmations and reminders.
    \item Test framework for unit/integration tests.
    \item CI runner and staging deployment target.
\end{itemize}
We will use these mature components to focus on workflow design, trust, and measurement rather than inventing new infrastructure.

\section{Project scope and deliverables}
\subsection{Initial deliverable}
\begin{itemize}[leftmargin=0.7cm]
    \item Traveler search + booking flow for nearby partners
    \item Partner onboarding + listing management (capacity, price, hours)
    \item QR-based drop-off and pickup handoff workflow
    \item Booking status tracking (Booked $\rightarrow$ Dropped $\rightarrow$ Stored $\rightarrow$ Picked up)
    \item Basic logging and timestamps for BFT measurement
    \item CI: build + backend unit tests + linting
    \item CD to staging with a single command or pipeline job
\end{itemize}

\subsection{Subsequent deliverables}
\begin{itemize}[leftmargin=0.7cm]
    \item Payments integration and automated partner payouts
    \item Ratings and reviews for partners
    \item Duplicate/quality checks for new partner listings
    \item Admin dashboard for coverage, utilization, and revenue metrics
    \item SMS/WhatsApp notifications for booking confirmations and reminders
    \item Improved geo-search, filtering, and indexed queries for scale readiness
\end{itemize}

\section{Risks and mitigations}
\begin{itemize}[leftmargin=0.7cm]
    \item \textbf{Trust and safety of stored items:} photo confirmation at drop-off/pickup, clear liability terms, verified partner onboarding, and an optional insurance add-on in a later iteration.
    \item \textbf{Partner adoption:} keep onboarding and day-to-day use extremely simple, minimize the steps needed to confirm a drop-off/pickup, and offer clear, transparent payout terms.
    \item \textbf{Payment/logistics risk:} hold payment until pickup is confirmed (escrow-style), rather than releasing it immediately at drop-off.
    \item \textbf{Evaluation measurement ambiguity:} instrument timestamps and define clear booking/handoff states before development begins.
    \item \textbf{Operational issues in pilot:} use a staging environment early and ensure CI/CD produces repeatable deployments.
\end{itemize}

\section{Summary}
This project proposes POTLI, a deployable web platform that connects
travelers who need short-term luggage storage with local shop and business
owners who have spare space to offer. It meets the course criteria by
emphasizing time-to-value through rapid iterations, implementing CI/CD to
support frequent safe releases, and using measurable evaluation criteria
(especially booking fulfillment time). It stays focused on engineering
workflow, trust, and scalability readiness rather than research-driven
novelty, while creating a genuine extra income opportunity for local
partners.

\end{document}